%=============================================================================!
% 9.0  CHAPTER OPENING                                                        !
%=============================================================================!
Once the forces are known, a simulation must use them to advance the
positions and velocities. The computer cannot follow every instant: it
computes a sequence of states separated by a finite time step. Chapter~2
used the forward Euler method, Chapter~3 used a more accurate solver, and
Chapter~7 showed that two methods with similar errors over one step can
behave quite differently over a long run.

This chapter develops velocity Verlet, a method widely used in molecular
dynamics. Its construction from Taylor series gives the error of a step;
its construction from kicks and drifts explains its reversibility and
preservation of phase-space volume. For a spring it keeps a slightly
different energy exactly, which lets us derive both its energy error and
its stability limit. These results give a practical way to choose the
step for a molecule, including the water example left open in
Chapters~1, 2 and~4. Long runs then test the method against RK4 and show
why the smooth forces of Chapter~8 matter. Chapter~10 assembles these
parts into a complete simulation in a periodic box.
